\subsection{Accuracy-Fairness Efficient Frontiers}
\label{sec:pareto}

We begin by examining the efficient frontier of accuracy vs. fairness for
the six datasets.  These curves are shown in Figure~\ref{fig:pareto},
and are obtained by varying the weight $\lambda$ on the fairness
regularizer, and for each value of $\lambda$ finding the model which
minimizes the associated regularized loss function.  For the logistic
regression cases, we extract probabilities from the learned model
$\vw$ as
$\Pr[y_i = 1] = {\exp(\vw \cdot x_i)}/{(1 + \exp(\vw \cdot x_i))}$ and
evaluate these probabilities as predictions for the binary labels
using $\mse$.  \footnote{Note that assessing the $\mse$ of these
  probabilities, interpreted as predictions, is a sensible choice.
  Since squared error is a proper scoring rule, if the labels are
  indeed generated according to a logistic regression model,
  minimizing the squared error of a predictor using mean squared error
  will elicit the true model as its minimizer.}
% \footnote{Note that assessing the
% accuracy loss via MSE is indeed a sensible choice: one assumes some
% true real-valued model applied to $\vx represents the probability of
% $\vx$'s binary observable  being . If one takes this stance,
% assessing the accuracy loss of a predictor using its $\mse will
% elicit the true model as its minimizer.}
In all of the datasets, as $\lambda$
increases, the models converge to the best constant predictor,
which minimizes the fairness penalties.
% which has $\mse$ of $1$ for regression datasets and $\mse$ of 0.25 for
% classification datasets.

Perhaps the most striking aspect of Figure~\ref{fig:pareto} is the
great diversity of tradeoffs across different datasets and different
fairness regularizers. For instance, if we examine the individual fairness
regularizer, on four of the datasets (Adult, Communities and Crime, Law School
and Sentencing), the curvature is relatively mild and constant --- there is
an approximately fixed rate at which fairness can be traded for accuracy.
In contrast, on COMPAS and Default, fairness loss can be reduced almost for ``free''
until some small threshold value, at which point the accuracy cost increases
dramatically. Similar comments can be made regarding hybrid fairness in the
logistic regression cases.

Individual fairness appears to be strictly more costly than group
fairness for the entire regime between the extremes of $\lambda = 0$
and $\lambda\rightarrow \infty$ for a majority of these datasets (with
the exception of COMPAS and Default datasets). In COMPAS and Default,
for small amounts of unfairness, group unfairness may be more costly
than the same level of individual unfairness.

Perhaps surprisingly, building separate models for each population
barely improves the tradeoff (and in some cases, hurts the
tradeoff for some values of $\lambda$) across almost all datasets and
fairness regularizers. This suggests that the academic discussion
about whether to allow disparate treatment (as explicitly allowed in
e.g. \cite{DworkHPRZ12, JosephKMR16}) -- i.e. whether sensitive
attributes such as race and gender should be ``forbidden'' vs. used in
building more accurate models for each subpopulation, is perhaps less
consequential than expected (at least on these datasets and using
linear or logistic regression models).

Note that in theory \grp is strictly less costly than 
\ind for any particular model $\vw$ (by
Jensen's inequality), and using separate models (one for each group)
should strictly improve the fairness/accuracy trade-off for any of
these notions of fairness. However, both \grp and separate models are more prone to overfitting 
(than \ind and single model), and hence a larger $\ell_2$-regularization parameter $\gamma$ 
tends to be selected in cross validation in these settings (see Appendix~\ref{sec:cross} for more details). 
This is a surprising interaction between the strength of the fairness penalty and the generalization ability 
of the model, and results in group fairness sometimes having a more severe tradeoff with accuracy 
when compared to individual fairness, and separate models having little benefit out of sample, although 
they can appear to have a large effect in-sample (because of the effects of over-fitting).

% Two models appears to help at
% most modestly improve the tradeoffs between any of our fairness
% notions and accuracy, for every dataset. Notable exceptions are \ind
% for the Pennsylvania Crime and Adult datasets
%\jm{Also for Adult? Someone is
%  checking on this, right?}
%  and slight improvements for two models \grp
% over one model \grp for the COMPAS and Adult
% datasets.

\iffalse
\begin{figure}[h!]
\begin{minipage}[b]{0.3\textwidth}
\centering
\includegraphics[width=4.5cm]{../../code/jamie/figures/diff-adult-no-log-unnormalized.jpg}
%\caption{}
\end{minipage}
\begin{minipage}[b]{0.3\textwidth}
\includegraphics[width=4.5cm]{../../code/jamie/figures/diff-communities-no-log-unnormalized.jpg}
%\caption{}
\end{minipage}
\begin{minipage}[b]{0.3\textwidth}
\includegraphics[width=4.5cm]{../../code/jamie/figures/diff-compas-no-log-unnormalized.jpg}
%\caption{}
\end{minipage}\\
\begin{minipage}[b]{0.3\textwidth}
\centering
\includegraphics[width=4.5cm]{../../code/jamie/figures/diff-default-no-log-unnormalized.jpg}
%\caption{}
\end{minipage}
\begin{minipage}[b]{0.3\textwidth}
\includegraphics[width=4.5cm]{../../code/jamie/figures/diff-law-no-log-unnormalized.jpg}
%\caption{}
\end{minipage}
\begin{minipage}[b]{0.3\textwidth}
\includegraphics[width=4.5cm]{../../code/jamie/figures/diff-philadelphia-no-log-unnormalized.jpg}
%\caption{}
\end{minipage}
\caption{Efficient frontiers of accuracy vs. fairness for each dataset. For datasets with binary-valued
targets (logistic regression), we consider three fairness notions (group, individual and hybrid), and
for each examine building a single model or separate models for each group, yielding a total of six curves.
For real-valued targets
(linear regression), we consider two fairness notions (group and individual), and again single or
separate models, yielding a total of four curves.
label{fig:pareto}}
\end{figure}

\fi

\iffalse
\jm{all below here old} \sj{The two figures show scaled
  fairness vs error for classification and regression datasets
  separately. Figures 1 and 2 are for individual fairness and figures
  3 and 4 are for group fairness. Some points to consider:
\begin{itemize}
\item as $\lambda\rightarrow \infty$ we converge to the best constant
  (and trivial) model. This model always predicts 0 for regression
  (which is the mean after normalization) and has MSE of 1. It
  predicts all positives/all negatives (whichever has a smaller error)
  for classification.
\item there is more variability if we start from a model that has a
  much better error rate than the trivial constant predictors
\end{itemize}
}
\fi
